\section{Experimental Details and Additional Results}
\label{sec:app-experiments}

\subsection{Implementation and correctness checks}
\label{sec:app-correctness}

\mname is implemented as a wrapper around the decoder stack that intercepts the
residual update of each block. Before running any experiment we verified the
wrapper against the following checks: (i) with no layers skipped, the wrapper
reproduces the original model's logits to within floating-point tolerance; (ii)
with $\alpha_\ell = 0$, \mname reproduces Plain Skip exactly; (iii) with
$\alpha_\ell = 1$, it reproduces Copy Update exactly; (iv) a manually skipped
layer agrees with the generic skip path; (v) residual updates are captured at the
block boundary of \cref{eq:residual}, after the block's own residual addition;
and (vi) padding tokens are excluded from the coefficient accumulators of
\cref{eq:fit}.

\subsection{Calibration and evaluation protocol}
\label{sec:app-protocol}

Coefficients are fitted on unlabeled WikiText-103 text \citep{merity2017pointer},
and $P_\ell$ is scored on a disjoint held-out split of the same size. The 0.5B
layer scan uses 16 calibration and 16 held-out sequences of length 512; the
larger models use \ph{32} of each. Statistics are accumulated in FP64
(\cref{eq:fit}) so that no hidden states need to be stored. Task evaluation uses
fixed example IDs with seed 42 throughout. Discovery and verification runs (0.5B, 1.5B) are executed in FP32
on an NVIDIA TITAN X (Pascal) GPU; the headline \topscale runs are executed in
BF16 on a single NVIDIA RTX 3090 GPU. Neither configuration uses FlashAttention,
so we use the reference attention implementation throughout. We caution that
absolute latencies are specific to this hardware, even though the relative
comparison between methods is not.

\subsection{Latency measurement}
\label{sec:app-latency}

We report \emph{prefill} latency only: a forward pass with
\texttt{use\_cache=False}, batch size \ph{8}, at sequence lengths 512 and 2048,
with \ph{10} warm-up iterations followed by \ph{30} measured iterations, and
\texttt{torch.cuda.synchronize()} around each timed region. We report the median
and interquartile range. We make no claim about end-to-end generation latency,
which depends on decode-time behaviour we do not measure here.

\subsection{Layer selection}
\label{sec:app-selection}

Unless stated otherwise, the skipped set is the $k$ layers of highest held-out
$P_\ell$ subject to a non-adjacency constraint: no two skipped blocks may be
consecutive, so that a skipped block never consumes a \emph{predicted} update as
the input to its own prediction. We compare this predictability-based selection
against a low-damage selection rule, which ranks layers by the NLL increase
caused by skipping each one individually. The comparison separates a failure of
update \emph{prediction} from a failure of layer \emph{selection}; because \mname
is applied at an identical layer set as every baseline in \Cref{tab:main}, the
main result is unaffected by which rule is used.

\subsection{\mname(2) and additional budgets}
\label{sec:app-ar2}

\Cref{tab:main} reports \mname(1), the single-scalar predictor of
\cref{eq:ar1}. We also evaluated the two-tap variant \mname(2) of
\cref{eq:ar2}, which fits $(\alpha_\ell, \beta_\ell)$ through a $2\times2$ ridge
system. We present \mname(1) as the primary method because its simplicity is part
of the contribution: a single scalar per skipped block is the smallest possible
departure from the zero-update assumption, and it is the variant whose cost is
most obviously negligible.
